

%-----------------------------------------------------
%
%    Contents
%
%-----------------------------------------------------

\sectionfix
\section{Implementation}
\label{sec:implementation}
\subsection{Work flow}
\begin{itemize}
\item GUI
\item Detailed step-by-step manual on how to use the application
\item tailored to the methods chapter mentioned above
\end{itemize}
\subsection{Dependencies}
A great benefit of Scansify is that it does not rely onto any major image processing library like openGL or PCL but rather makes use of the native implementation framework KinectFusion \cite{Izadi:2011:KRR:2047196.2047270}. 
\subsubsection{Kinect Fusion}
This framework allows users to create high-quality 3D reconstructions of indoor rooms in real-time by moving the Kinect around the scene. The core system provides a great amount of utilities which we uses to achieve our goal. Unfortunately, they only provide us with an api structure which imposes a restriction on the use of their source code. Partiularly, our pipeline extends their implementation of a shader and surface reconstruction.
The function of these components and how we integrated them into our project will be enlighten in more detail in the respective paragraph later on.
\subsection{Kinect Fusion}
\begin{itemize}
\item Introduction
\begin{itemize}
\item system allows a user to pickup a standard Kinect camera and move rapidly within a room to reconstruct a high-quality, geometrically precise 3D model of the scene
\item avoids the reliance on RGB allowing use in in- door spaces with variable lighting conditions
\item reconstructing surfaces, which more accurately approximate real-world geometry
\item aim to perform camera track- ing without the need for prior augmentation of the space
\end{itemize}
\item What do we use
\item remark regarding restricted source code access
\item tailored approach
\end{itemize}
\subsection{Software architecture}
The software pipeline described in the previous chapter (see \autoref{subsec:overview}) was implemented as a Microsoft Foundation Classes ( MFC ) application and written in C++. MFC is a collection of frameworks created by Microsoft for programming systems with graphical user interfaces \cite{Fayad:1997:OAF:262793.262798}. The purpose of the different modules are explicitely explained in \autoref{sec:goals}. Therefore, following exempt will rather give an outlook on how these components interact with each other and how they were realized.
\subsubsection{Data processing}
The depth stream is collected via the Kinect software development kit ( SDK ) which is the official programming library for processing raw input from the Kinect \cite{Abhijit:2012:KWS:2490784}. Hereby, data is presented as numerical values that needs to be mapped to the real world coordinates. Since the point cloud scan is quite noisy initially, a conditional and outlier filter technique is applied to improve more smooth and accurate data aquisition. These methods essentially remove extreme points that are behaving extraordinary due to their position in relation to their neighbours.\\
Afterwards, multiple point clouds are merged together by using advanced volumetric representation \cite{Curless:1996:VMB:237170.237269} and ICP which is provided by Kinect Fusion.
\subsubsection{Rendering}
A major challenge was the optimization of our rendering pipeline which is entirely responsible for displaying the digital model onto our GUI. Herewith, we implemented a concurrent raytracer that traces a certain direction from a perspective camera angle and then labels all 3D points from that view. At the end, each image pixel will have computed their own collision point with the 3D environment. Such scenes often consist of thousands of polygons and therefore needs to be incapsulated into accurate bounding boxes to pre-eliminate candiates that are obviously not in the area of the current shooting ray. In our case, we realized the slabs-based technique which improves processing speed drastically by reducing the amount of intersection calculations \cite{mensmann2010slab}.\\
 To further enhance performance, an implementation of the bounding volume hierarchy ( BVH ) \cite{Karras:2013:FPC:2492045.2492055} algorithm and surface area heuristic ( SAH ) \cite{Karras:2013:FPC:2492045.2492055} are integrated. These approaches provide an efficient way to build up an acceleration structure for traversing through the nodes inside of the BVH tree. On a side note, polygons are usually presented as edged triangles but in our case we use a smooting method to elevate a more accurate geometrical representation ( see comparision...).
Afterwards, we proceed to use the optimized shader from KinectFusion to map each digitial 3D point to a physically realistic color according to the ambient light environment. The image renderer utilizes the GPU for multi-threading processing and is hence able to sample in real-time.
\subsubsection{Mouse picking and camera manipulation}
The user is able to draw on the digital model by using the mouse and clicking on the target area. As soon as an mouse signal passes on, our system creates a new camera instance with the current position of the cursor. Afterwards, a ray is shot through that view to determine which exact area is being annotated by the user (drawing). With that information, Scansify is able to store the draw input inside of the 3D model and update the view accordingly.\\
In addition, the camera view inside of the digital environment can also be steered, translated and rotated. We utilize an arcball camera implementation to achieve that goal \cite{Shoemake:1992:AUI:155294.155312}.
\raggedbottom
\subsubsection{Unwrapping a 3D volume}
\begin{itemize}
\item why do we want this
\item ongoing HCI challenges regarding this domain
\item Math behind it with pseudocode
\end{itemize}

Once the user starts drawing onto the model, Scansify slowly fetches the input and builds an interconnected graph with it. As the graph expands, 

\subsubsection*{Integrating seams}
How does using seams work with our algorithm.
